\chapter{Hash}
\section{Dizionari}
Una struttura dati molto comoda ed 
efficiente per supportare diverse applicazioi è quelle del \textbf{dizionario}
\begin{itemize}
    \item consente di rappresentare il concetto matematico di relazione univoca R : D $\rightarrow$ C, o associazione chiave valore
    \item deve supportare in modo efficiente le operazioni: Insert, Search, Delete  
\end{itemize}
\noindent Esistono diversi modi per realizzare un dizionario tra cui:
\begin{itemize}
    \item tebelle ad indirizzamento diretto
    \item tabelle di hash
\end{itemize}
\section{Tabelle ad indirizzamento diretto}
Il dizionario si gestisce in maniera semplice sfruttando. un 'array di puntatori'\\
Dato l'insieme delle chiavi $U=(k_1,k_2,..,k_m)$ possiamo usare un array T in cui associo ad ogni posizione una chiave $k_i$\\
$T[k]$ è un puntatore al valore associato alla chiave k.
AGGIUNGI FOTO\\
Le operazioni caratteristiche della struttura dati possono essere realizzate in modo semplice\\ 
AGGIUNGI FOTO\\
Pro:\\
\begin{itemize}
    \item Operazioni efficienti eseguite in $O(1)$
\end{itemize}
\noindent Contro:
\begin{itemize}
    \item occorre riservare memoria sufficiente per tante celle quante sono le possibili chiavi 
    \item Se U è molto grande è inutilizzabile a causa delle limitazioni di memoria
    \item Se le chiavi effettive sono molto di meno delle chiavi potenziali abbiamo uno spreco di memoria
\end{itemize}

\section{Tabelle di Hahs}
Concettualmente si realizza con gli stessi principi utilizzati per le tabelle di 
indirizzamento diretto ma l'associazione chiave-cella è definta tramite una funzione e non in maniera statica\\
In una \textbf{tabella hash} di m celle ogni chiave k viene memorizzata nella cella $h(k)$ usando una funzione $h: U \rightarrow {0,..,m-1}$ detta \textbf{funzione di hash}\\
Vantaggio: richeide memoria proporzionale al numero massimo di chiavi presenti nel dizionario indipendentemente dalla cardinalità dell'insieme $U$ di tutte le possibili chiavi.\\
AGGIUNGI FOTO\\
\subsection{Collisioni}
\textbf{Il principio della PICCIONAIA}
\begin{definition}
    Se p piccioni devono trovare posto in c caselle e ci sono più piccioni che caselle allora
    in qualche casella entreranno almeno due piccioni
\end{definition}
\noindent Siccome $|U|>m$, esisteranno coppie di chiavi diverse $k_1 \neq k_2$ tali che $h(k_1)=h(k_2)$ diremo in questo caso che vi è una coliisione tra le due chiavi

\subsection{Risoluzione collisioni}
\subsection{Metodo del concatenamento}
L'insieme di input di $n$ viene diviso casualmente in $m$ sottoinsiemi, una funzione di hash $h()$ determina a quale sottoinsieme appartiene un elemento, ogni. sottoinsieme è gestito in modo indipendente tramite una lista\\
AGGIUNGI FOTO\\
AGGIUNGI FOTO o MINIPAGE A 3 REALIZZAZIONE\\
\subsection{Analisi dell'hashing con concatenamento}
Caso peggiore\\
Tutte le $n$ chiavi sono in conflitto e puntano tutte alla stessa cella\\
AGGIUNGI FOTO\\

\subsection{Progettare funzioni di hash}
AGGIUNGI FOTO\\
Una buona funzione di hash dovrebbe soddisfare l'ipotesi di hash uniforme indipendente\\
Ogni chiave. ha la stessa probabilità $1/m$ di essere mandata di ua qualsiasi delle $m$ celle, indipendentementedalle chiavi inserite precedentemente\\
Svantaggio:\\
L'ipotesi di hash unifore indipendente dipende dalle probabilità con cui vengono estratti gli elementi da inserire\\
Assunzione:\\
Le chiavi che considereremo nelle prossime pagine sono rappresentate degli interi non negativo piccoli
\subsection{Hashing Statico}
\textbf{Metodo divisione}\\
Associa ad una chiave $k$ una delle $m$ celle prendendo il resto della divisioe intera tra $k$ e $m$\\
AGGIUNGI FOTO\\
\textbf{Meotdo della moltiplicazione}
Moltiplica la chiave k per una costante A (A comoresa tra 0 e 1) ed estrai la parte frazionaria di kA\\
Moltiplica il valore ottenuto per m e prendi la parte intera inferiore\\
$$h(k)=\lfloor m(kA \mod 1) \rfloor$$\\
AGGIUNGI FOTO\\
\textbf{Metodo della moltiplicazione}\\
Mtolo efficiente se m è una potenza di 2\\
$m=2^p$ e $A=q/2^W$\\
Assumiamo che ogni chiave possa essere rappresentata da una cola parola di w bit\\
FOTO
\subsection{Warning}
Nessuna funzione di hash può evitare che un avversario malevolo inserisca
nella tabella una sequenza di valori che vadano a finire tutti nella stessa lista\\
Più in generale, per ogni funzione di hash si possono trovare delle distribuzioni
di probabilità degli input per le quali la funzione non ripartisce bene le chiavi
tra le varie liste della tabella hash
\subsection{Hashing Randomizzato}
Possiamo usare la randomizzazione per rendere il comportamento della tabella hash indipendente dall'input.\\
L’idea è quella di usare una funzione di hash scelta casualmente in un insieme
“universale” di funzioni di hash.
Questo approccio viene detto \textbf{hashing universale}.\\
Un insieme $H$ di funzioni di hash che mandano un insieme U di chiavi bell'insieme {0,1,..,m-1} degli indici della tabella hash si dice universale se per ogni coppia di chiavi diverse $j$ e $k$ vi sono al più $|H|/m$ funzioni di hash in $H$ tali che $h(j)=h(k)$\\
continua pagina 21 e 22
\subsection{Risoluzione delle collisioni}
\textbf{Metodo dell'indirizzamento aperto}\\
Con la tecnica di indirizzamento aperto tutti gli elementi stanno nella tabella, non utilizziamo puntatori a zone estere di memoria ma tutto risiede nella tabella,
se la tabella è piena potrei non riuscire a completare un inserimento\\
Quando voglio inserire un elemento, provo a mattelro nella posizione primaria se non riesco, provo a inserirlo nella posizione secondaria, se non riesco nella terziaria e cosi via...\\
Ne segue che:
La funzione di hash non individua una singola cella ma ma deve definire un ordine in cui ispezionare
tutte le celle. \\
L’inserimento di un elemento avviene nella prima cella libera che si incontra nell’ordine di ispezione.\\
In questo caso, la funzione di hash è una funzione h(k,i) che al variare di $i$ tra $0$ ed $m-1$ fornisce, per ciascuna chiave k, una sequenza di indici\\
$$h(k,0), h(k,1),..., h(k,m-1)$$\\
che rappresenta l'ordine di ispezione\\
Siccome vogliamo poter ispezionare tutte le celle, la sequenza deve essere una permutazione dell'insieme degli indici $0,1,..,m-1$ della tabella.\\
FOTO\\
\textbf{Metodo dell'indirizzamento aperto-delete}\\
\textbf{warning}
Non possiamo infatti limitarci a porre nil (null) nella cella, perché altrimenti
interromperemmo delle sequenze di ispezione\\
\textbf{Delete} assegna alla chiave dell'elemento da togliere un particolare diverso da ogni possibile chiave:\\
FOTO\\
Con l'indirizzamento aperto, la funzione di hash fornisce una sequenza di ispezione\\
Si modifica l’ipotesi di hashing indipendente e uniforme e si parla di hashing con
permutazioni indipendenti e uniformi (più comunemente noto come hashing
uniforme)\\
ogni chiave ha la stessa probabilità $1/m!$ di generare
una qualsiasi delle $m!$ possibili sequenze di ispezione\\
Ci sono tre tecniche comunemente usate per determinare l’ordine di ispezione:\\
\begin{enumerate}
    \item Ispezione lineare
    \item Ispezione quadratica
    \item Doppio hashing
\end{enumerate}
\noindent Nessuna delle tre genera tutte le $m!$ sequenze di ispezione. Le prime due ne generano soltanto m e l’ultima ne genera $m^2$\\
Ispeazione lineare\\
La funzione di hash $h(k,i)$ si ottiene da una funzione di hash ordinaria $h'(k)$
ponendo\\
$$ h(k,i) = [h'(k)+i] \mod m $$\\
L’esplorazione inizia dalla cella $h(k,0) = h'(k)$ e continua con le celle $h'(k)+1$,
$h'(k)+2$, ecc. fino ad arrivare alla cella $m-1$, dopo di che si continua con le celle
$0,1,ecc..$ fino ad aver percorso circolarmente tutta la tabella.\\
L’ispezione lineare è facile da implementare ma soffre del problema del
clustering primario:\\
i nuovi elementi inseriti nella tabella tendono ad addensarsi (formare un
cluster) attorno agli elementi già presenti\\
Ispezione quadratica\\
La funzione di hash $h(k, i)$ si ottiene da una funzione di hash ordinaria $h'(k)$
ponendo\\
$$ h(k,i)=[h'(k)+c_1 i+ c_2 i^2] \mod m$$\\
dove $c_1$ e $c_2$ sono due costanti con $c_2 \neq 0$\\
I valori di $m$, $c1$ e $c2$ non possono essere qualsiasi ma debbono essere scelti
opportunamente in modo che la sequenza di ispezione percorra tutta la tabella\\
ispezione doppio hashing\\
La funzione di hash $h(k,i)$ si ottiene da due funzioni di hash ordinarie $h_1(k)$ ed 
$h_2(k)$ ponendo\\
$$ h(k,1)=[h_1(k)+ih_2 (k)]\mod m $$\\
Perché la sequenza di ispezione percorra tutta la tabella, il valore di $h_2 (k)$ deve
essere relativamente primo con m.\\
Possiamo soddisfare questa condizione in diversi modi.\\
FOTO\\
\section{Analisi della Complessità}
\textbf{Caso Peggiore}\\
Tutte le n chiavi sono in conflitto e puntano tutte alla stessa cella.\\
FOTO\\
\textbf{caso Medio}\\
Le prestazioni dipendono dal modo in cui la funzione di hash distribuisce 
le chiavi tra le celle.\\
Valutiamo la complessità media di Search in funzione del \textbf{fattore di carico} $a = n/m$.\\
Valutiamo le prestazioni sotto l’ipotesi che la funzione hash $h(k)$ sia uniforme e
indipendente\\
i.e., $h(k)$ distribuisce in modo uniforme le n chiavi tra le m liste.
\begin{enumerate}
    \item ogni elemento in input ha la stessa probabilità di essere mandato in una qualsiasi delle m celle
    \item l’immagine tramite la funzione di hash di ogni elemento è indipendente da quella di ogni altro elemento
\end{enumerate}
\noindent FOTO PAGINA 39\\
\textbf{Teorema 1}\\
\begin{definition}
In una tabella hash le cui collisioni sono risolte tramite la tecnica del concatenamento, una
ricerca senza successo richiede in media un tempo $\Theta(1 + \alpha)$, nell’ipotesi di hashing
uniforme e indipendente.
\end{definition}
\noindent Dimostrazione:\\
Consideriamo che il calcolo della funzione $h(k)$ sia O(1)\\
\textbf{Search}\\
Calcola $j=h(k)$\\
Controlla tutti gli $n_j$ elementi della lista $T[j]$\\
Nell'ipotesi di hash uniforme indipendente $E[n_j]=\alpha$ e quindi l'algoritmo richiede un tempo $\Theta(1+\alpha)$ in media.\\
\textbf{Teorema2}
\begin{definition}
In una tabella hash le cui collisioni sono risolte tramite la tecnica del concatenamento, una
ricerca con successo richiede in media un tempo $\Theta(1 + \alpha)$, nell’ipotesi di hashing
uniforme e indipendente.
\end{definition}
\noindent dimostrazione\\
CONTINUA DA PAGINA 42
